% TCC: História da Engenharia de Reservatórios
\PassOptionsToPackage{brazilian}{babel}
\documentclass[12pt,a4paper,oneside,brazilian,sumario=tradicional]{abntex2}

% CORREÇÃO DO IDIOMA (evita aviso do babel)
\PassOptionsToPackage{brazilian}{babel}

\usepackage[T1]{fontenc}
\usepackage[utf8]{inputenc}
\usepackage{times}
\usepackage{geometry}
\geometry{a4paper,left=3cm,right=2cm,top=3cm,bottom=2cm}
\usepackage{setspace}
\OnehalfSpacing
\usepackage{indentfirst}
\setlength{\parindent}{1.25cm}
\usepackage{emptypage}

% Figuras, tabelas e ilustrações
\usepackage{graphicx}
\usepackage{booktabs}
\usepackage{longtable}
\usepackage{array}
\usepackage{multirow}
\usepackage{float}
\usepackage{caption}
\usepackage{tikz}
\usetikzlibrary{positioning, arrows.meta, shapes, patterns, decorations.pathmorphing, calc}
\captionsetup{font=small, labelfont=bf, justification=centering, skip=6pt}
\renewcommand{\fonte}[1]{\vspace{2pt}\par\noindent\footnotesize Fonte: #1}

\usepackage{amsmath}
\usepackage{amssymb}
\usepackage{amsfonts}

\usepackage{hyperref}
\usepackage{xcolor}
\usepackage{enumitem}
\usepackage{lastpage}
\usepackage{epigraph}
\usepackage{microtype}

\usepackage[alf,abnt-etal-list=0,abnt-etal-cite=3,abnt-repeated-author-omit=yes,abnt-emphasize=bf]{abntex2cite}

\hypersetup{
  pdftitle={História da Engenharia de Reservatórios},
  pdfauthor={Nome do Estudante},
  colorlinks=true,
  linkcolor=black,
  citecolor=black,
  urlcolor=black
}

% Compatibilidade para entradas geradas no .bbl (comandos antigos)
\DeclareOldFontCommand{\tt}{\normalfont\ttfamily}{\texttt}
\usepackage{url}
\providecommand{\htmladdnormallink}[2]{\href{#2}{#1}}

% Dados da capa
\titulo{História da Engenharia de Reservatórios}
\autor{Nome do Estudante}
\orientador{Nome do Orientador}
\instituicao{%
  Instituto Superior Politécnico de Tecnologias e Ciências -- ISPTEC\\
  Departamento de Geociências\\
  Curso de Engenharia de Petróleo
}
\local{Luanda -- Angola}
\data{2026}
\tipotrabalho{Trabalho de Avaliação Contínua}
\preambulo{Trabalho investigativo apresentado ao Curso de Engenharia de
Petróleo do Departamento de Geociências do Instituto Superior
Politécnico de Tecnologias e Ciências (ISPTEC) como avaliação contínua
na disciplina de Engenharia de Reservatórios, sob orientação do
\imprimirorientador.}

\begin{document}
\frenchspacing
\pretextual
\imprimircapa
\imprimirfolhaderosto

\begin{fichacatalografica}
  \thispagestyle{empty}
  \vspace*{\fill}
  \begin{center}
    \fontsize{9}{9}\selectfont
    \noindent
    \imprimirautor.

    \imprimirtitulo / \imprimirautor. -- \imprimirlocal: ISPTEC, \imprimirdata.

    \pageref{LastPage} p. : il.

    Trabalho de Avaliação Contínua -- Instituto Superior Politécnico de
    Tecnologias e Ciências (ISPTEC). Departamento de Geociências. Curso
    de Engenharia de Petróleo, \imprimirlocal, \imprimirdata.

    Professor: \imprimirorientador.

    1. Engenharia de Reservatórios -- História. 2. Petróleo -- Produção.
    3. Simulação Numérica. 4. Recuperação Avançada de Petróleo (EOR).

    I. \imprimirorientador. II. Instituto Superior Politécnico de Tecnologias e Ciências.
  \end{center}
\end{fichacatalografica}

\begin{folhadeaprovacao}
  \thispagestyle{empty}
  \begin{center}
    {\ABNTEXchapterfont\large\MakeUppercase{\imprimirautor}}
    \vspace*{1cm}
    \begin{center}
      \ABNTEXchapterfont\bfseries\large\MakeUppercase{\imprimirtitulo}
    \end{center}
    \vspace*{1cm}
    \hspace{.38\textwidth}
    \begin{minipage}{.56\textwidth}
      \imprimirpreambulo
    \end{minipage}
    \vspace*{1cm}
  \end{center}

  Trabalho avaliado em \underline{\hspace{1.8cm}} de
  \underline{\hspace{3.5cm}} de \underline{\hspace{1.2cm}}

  \vspace{1.5cm}
  \assinatura{\textbf{\imprimirorientador} \\ Prof. Orientador}

  \begin{center}
    \imprimirlocal, \imprimirdata
  \end{center}
\end{folhadeaprovacao}

\setlength{\absparsep}{18pt}
\begin{resumo}
Este trabalho apresenta uma visão histórica e crítica da Engenharia de
Reservatórios, desde as práticas empíricas iniciais até as técnicas
contemporâneas de simulação, recuperação avançada (EOR) e desafios
ambientais. O texto sintetiza marcos históricos, métodos de estimativa de
volumes, principais mecanismos de produção, e perspectivas futuras com
ênfase na transformação digital e sustentabilidade. O documento foi
construído a partir de material de referência e fontes técnicas
especializadas, obedecendo normas acadêmicas para apresentações em LaTeX.
\vspace{6pt}
\noindent\textbf{Palavras-chave}: Engenharia de Reservatórios; EOR; simulação; história; Angola.
\end{resumo}

\begin{resumo}[Abstract]
This study provides a concise historical overview of Reservoir Engineering,
from early empirical practices to contemporary techniques including
numerical simulation and enhanced oil recovery (EOR). It highlights the main
historical milestones, recovery methods and future perspectives, emphasizing
digital transformation and sustainability.
\end{resumo}

% Listas pré-textuais conforme modelo ISPTEC
\pdfbookmark[0]{\listfigurename}{lof}
\listoffigures*
\cleardoublepage

\pdfbookmark[0]{\listtablename}{lot}
\listoftables*
\cleardoublepage

\begin{siglas}
  \item[CCS]    \textit{Carbon Capture and Storage} -- Captura e Armazenamento de CO\textsubscript{2}
  \item[EOR]    \textit{Enhanced Oil Recovery} -- Recuperação Avançada de Petróleo
  \item[IA]     Inteligência Artificial
  \item[IoT]    \textit{Internet of Things} -- Internet das Coisas
  \item[ISPTEC] Instituto Superior Politécnico de Tecnologias e Ciências
  \item[OOIP]   \textit{Original Oil In Place} -- Volume Original de Óleo no Reservatório
  \item[PVT]    Análise Pressão-Volume-Temperatura
\end{siglas}

\begin{simbolos}
  \item[$q$] Vazão volumétrica (cm\textsuperscript{3}/s)
  \item[$k$] Permeabilidade absoluta (Darcy ou mD)
  \item[$A$] Área da seção transversal (cm\textsuperscript{2})
  \item[$\mu$] Viscosidade dinâmica do fluido (cP)
  \item[$\phi$] Porosidade efetiva (fração)
  \item[$N$] Volume original de óleo \textit{in place} (OOIP)
\end{simbolos}

\pdfbookmark[0]{\contentsname}{toc}
\tableofcontents*
\cleardoublepage

\textual

\chapter{Introdução}
A engenharia de reservatórios constitui uma área essencial na indústria do
petróleo, dedicada à análise, gestão e otimização da produção de
hidrocarbonetos em meios porosos. Sua origem está ligada ao desenvolvimento
da indústria petrolífera comercial a partir de meados do século XIX e ao
avanço de conceitos físicos que permitiram a modelagem do escoamento em
subsuperfície \cite{yergin1991,dake1994}.

Este trabalho organiza conteúdo histórico e técnico para formar um TCC
compatível com as normas acadêmicas, incluindo capa, resumo, sumário,
introdução, desenvolvimento, conclusão e referências, empregando o pacote
`abntex2` para conformidade com as normas ABNT.

\chapter{Objetivos}
\section{Objetivo Geral}
Analisar a evolução histórica da Engenharia de Reservatórios e sua
importância técnica e institucional para a gestão eficiente de campos
petrolíferos.

\section{Objetivos Específicos}
\begin{itemize}
  \item Descrever a origem da indústria petrolífera e a emergência da engenharia de reservatórios;
  \item Apresentar a evolução dos métodos de recuperação (primária, secundária e EOR);
  \item Identificar desafios atuais e perspectivas futuras, com foco em Angola.
\end{itemize}

\chapter{Metodologia}
Trata-se de uma pesquisa bibliográfica e qualitativa, com levantamento e
análise crítica de fontes técnicas, livros e publicações setoriais.
As informações foram sintetizadas para atender à finalidade educativa e
acadêmica do TCC \cite{ahmed2010,craft1991}.

\chapter{Desenvolvimento}
\section{Origem da indústria petrolífera}
O marco inicial da indústria petrolífera moderna é comumente atribuído à
perfuração do primeiro poço comercial por Edwin Drake, em 1859, nos EUA;
posteriormente, a indústria expandiu-se globalmente, exigindo técnicas de
exploração e produção cada vez mais sofisticadas \cite{yergin1991}.

\section{Evolução histórica da Engenharia de Reservatórios}
A disciplina evoluiu de práticas empíricas para uma área científica com base
em princípios físicos como a Lei de Darcy, metodologias de balanço de
materiais e, mais tarde, simulação numérica \cite{muskat1949,dake1994,peaceman1977}.
Com o advento da computação, modelos tridimensionais e a integração de
dados geológicos e de produção tornaram-se instrumentos essenciais
para o planejamento de produção \cite{peaceman1977,aziz1979}.

\section{Métodos de recuperação de petróleo}
Os métodos de recuperação evoluíram de mecanismos naturais (recuperação
primária) a técnicas secundárias (injeção de água/gás) e, ultimamente, a
EOR (injeção de CO$_2$, polímeros, surfactantes, injeção de vapor) para
recuperar volumes residuais \cite{lake2007,craft1991}.

\section{Desafios atuais e perspectivas}
Desafios modernos incluem a gestão de campos maduros, complexidade
geológica, exigências ambientais e volatilidade de preços. A integração de
inteligência artificial e digital twins oferece novas oportunidades para a
otimização, enquanto a captura e armazenamento de carbono (CCS) e outras
tecnologias seguem como áreas prioritárias de desenvolvimento \cite{mungan2019,bachu2000}.

\chapter{Conclusão}
Este trabalho sintetiza a trajetória da Engenharia de Reservatórios, mostrando
como a disciplina transicionou de conhecimentos empíricos para uma prática
multidisciplinar e tecnológica. Para Angola, recomenda-se priorizar a
capacitação local, investimento em dados e pilotos bem delineados para
aplicar métodos de EOR e considerar projetos de CCS quando viáveis.

\postextual
% Incluir todas as entradas do arquivo .bib (mesmo que não citadas diretamente)
\nocite{*}
\bibliographystyle{abntex2-alf}
\bibliography{referencias}

\end{document}
